% ═══════════════════════════════════════════════════════════════
% Figure: L0-L4 Adaptive Attack Model with Intelligence Chaining
% Usage: % ═══════════════════════════════════════════════════════════════
% Figure: L0-L4 Adaptive Attack Model with Intelligence Chaining
% Usage: % ═══════════════════════════════════════════════════════════════
% Figure: L0-L4 Adaptive Attack Model with Intelligence Chaining
% Usage: % ═══════════════════════════════════════════════════════════════
% Figure: L0-L4 Adaptive Attack Model with Intelligence Chaining
% Usage: \input{figures/fig_attack_model.tex}
% ═══════════════════════════════════════════════════════════════
\begin{figure}[t]
\centering
\resizebox{\columnwidth}{!}{%
\begin{tikzpicture}[
    >=stealth,
    level/.style={
        rectangle, rounded corners=6pt, draw=#1!70, fill=#1!10,
        line width=1pt, minimum width=3.0cm, minimum height=1.2cm,
        align=center, font=\sffamily\bfseries\small
    },
    intel/.style={
        rectangle, rounded corners=3pt, draw=black!40, fill=yellow!8,
        font=\sffamily\scriptsize, align=left, text width=3.5cm,
        inner sep=3pt
    },
    tool/.style={
        ellipse, draw=gray!60, fill=gray!8,
        font=\sffamily\tiny, inner sep=2pt, align=center
    },
    chain/.style={->, line width=1.8pt, draw=black!60},
    intflow/.style={->, line width=0.8pt, draw=orange!60, dashed},
    ttp/.style={
        rectangle, rounded corners=2pt, draw=teal!50, fill=teal!5,
        font=\ttfamily\tiny, inner sep=2pt
    },
]

% ── Level nodes (vertical cascade) ───────────────────────────
\node[level=gray]    (L0) at (0, 8)   {L0\\[-2pt]{\footnotesize Script Kiddie}};
\node[level=blue]    (L1) at (0, 6)   {L1\\[-2pt]{\footnotesize Basic}};
\node[level=yellow]  (L2) at (0, 4)   {L2\\[-2pt]{\footnotesize Intermediate}};
\node[level=purple]  (L3) at (0, 2)   {L3\\[-2pt]{\footnotesize Advanced}};
\node[level=red]     (L4) at (0, 0)   {L4\\[-2pt]{\footnotesize APT}};

% ── Cascade arrows ───────────────────────────────────────────
\draw[chain] (L0) -- (L1);
\draw[chain] (L1) -- (L2);
\draw[chain] (L2) -- (L3);
\draw[chain] (L3) -- (L4);

% ── Skill parameter (s) ─────────────────────────────────────
\node[font=\sffamily\tiny\bfseries, text=gray!60] at (-1.9, 8) {$s=0.20$};
\node[font=\sffamily\tiny\bfseries, text=blue!60] at (-1.9, 6) {$s=0.35$};
\node[font=\sffamily\tiny\bfseries, text=yellow!70!black] at (-1.9, 4) {$s=0.55$};
\node[font=\sffamily\tiny\bfseries, text=purple!60] at (-1.9, 2) {$s=0.75$};
\node[font=\sffamily\tiny\bfseries, text=red!60] at (-1.9, 0) {$s=1.00$};

% ── Attack actions (right side) ──────────────────────────────
\node[intel] (A0) at (4.5, 8) {
    \textbullet\ TCP SYN scan (20 ports)\\
    \textbullet\ UDP MAVLink probe\\
    \textbullet\ Banner grabbing
};
\node[intel] (A1) at (4.5, 6) {
    \textbullet\ HEARTBEAT $\to$ sysid 추출\\
    \textbullet\ PARAM\_REQUEST\_LIST\\
    \textbullet\ ARM/TAKEOFF 명령 주입\\
    \textbullet\ GPS spoof 시도
};
\node[intel] (A2) at (4.5, 4) {
    \textbullet\ HTTP API 열거\\
    \textbullet\ 브루트포스 로그인\\
    \textbullet\ 크레덴셜 재사용\\
    \textbullet\ 미션 데이터 추출
};
\node[intel] (A3) at (4.5, 2) {
    \textbullet\ CVE-2026-25253 exploit\\
    \textbullet\ RTSP DoS\\
    \textbullet\ 브레드크럼 체인 추적\\
    \textbullet\ WS skill\_invoke
};
\node[intel] (A4) at (4.5, 0) {
    \textbullet\ SSH lateral movement\\
    \textbullet\ Cross-drone cred reuse\\
    \textbullet\ 펌웨어 업로드\\
    \textbullet\ C2 persistence
};

\draw[->, gray!40] (L0.east) -- (A0.west);
\draw[->, gray!40] (L1.east) -- (A1.west);
\draw[->, gray!40] (L2.east) -- (A2.west);
\draw[->, gray!40] (L3.east) -- (A3.west);
\draw[->, gray!40] (L4.east) -- (A4.west);

% ── Intelligence harvest (far right) ────────────────────────
\node[ttp] (I0) at (9.0, 8) {open\_ports, banners};
\node[ttp] (I1) at (9.0, 6) {sysid, params, arm\_ok};
\node[ttp] (I2) at (9.0, 4) {api\_token, ssh\_pw, signing\_key};
\node[ttp] (I3) at (9.0, 2) {ws\_token, permissions};

\draw[->, gray!30] (A0.east) -- node[above, font=\tiny\sffamily, text=orange!70] {harvest} (I0.west);
\draw[->, gray!30] (A1.east) -- node[above, font=\tiny\sffamily, text=orange!70] {harvest} (I1.west);
\draw[->, gray!30] (A2.east) -- node[above, font=\tiny\sffamily, text=orange!70] {harvest} (I2.west);
\draw[->, gray!30] (A3.east) -- node[above, font=\tiny\sffamily, text=orange!70] {harvest} (I3.west);

% ── Intel chaining arrows (diagonal) ────────────────────────
\draw[intflow] (I0.south) -- (A1.east);
\draw[intflow] (I1.south) -- (A2.east);
\draw[intflow] (I2.south) -- (A3.east);
\draw[intflow] (I3.south) -- (A4.east);

% ── Intel chain label ────────────────────────────────────────
\node[rotate=-90, font=\sffamily\tiny\bfseries, text=orange!60] at (10.3, 4) {Intelligence Chaining};

% ── ATT\&CK TTPs (left side) ────────────────────────────────
\node[ttp] at (-3.2, 8) {T0846 Remote Svc};
\node[ttp] at (-3.2, 6) {T0836 Modify Param};
\node[ttp] at (-3.2, 4) {T0812 Default Cred};
\node[ttp] at (-3.2, 2) {T0821 Modify Task};
\node[ttp] at (-3.2, 0) {T0839 Module FW};

\draw[->, teal!30] (-2.0, 8) -- (L0.west);
\draw[->, teal!30] (-2.0, 6) -- (L1.west);
\draw[->, teal!30] (-2.0, 4) -- (L2.west);
\draw[->, teal!30] (-2.0, 2) -- (L3.west);
\draw[->, teal!30] (-2.0, 0) -- (L4.west);

% ── Protocol labels ──────────────────────────────────────────
\node[font=\sffamily\tiny, fill=white, draw=gray!30, rounded corners=2pt, inner sep=1.5pt] at (2.2, 8) {TCP/UDP};
\node[font=\sffamily\tiny, fill=white, draw=blue!30, rounded corners=2pt, inner sep=1.5pt] at (2.2, 6) {MAVLink};
\node[font=\sffamily\tiny, fill=white, draw=yellow!50, rounded corners=2pt, inner sep=1.5pt] at (2.2, 4) {HTTP};
\node[font=\sffamily\tiny, fill=white, draw=purple!30, rounded corners=2pt, inner sep=1.5pt] at (2.2, 2) {WebSocket};
\node[font=\sffamily\tiny, fill=white, draw=red!30, rounded corners=2pt, inner sep=1.5pt] at (2.2, 0) {SSH/Multi};

\end{tikzpicture}
}
\caption{L0--L4 적응형 공격 모델. 각 레벨은 이전 레벨에서 수확한 인텔리전스를 체이닝하여 다음 레벨의 공격에 활용한다. 오른쪽 점선은 인텔리전스 흐름(Intelligence Chaining)을, 왼쪽은 MITRE ATT\&CK for ICS TTP 매핑을 나타낸다.}
\label{fig:attack}
\end{figure}

% ═══════════════════════════════════════════════════════════════
\begin{figure}[t]
\centering
\resizebox{\columnwidth}{!}{%
\begin{tikzpicture}[
    >=stealth,
    level/.style={
        rectangle, rounded corners=6pt, draw=#1!70, fill=#1!10,
        line width=1pt, minimum width=3.0cm, minimum height=1.2cm,
        align=center, font=\sffamily\bfseries\small
    },
    intel/.style={
        rectangle, rounded corners=3pt, draw=black!40, fill=yellow!8,
        font=\sffamily\scriptsize, align=left, text width=3.5cm,
        inner sep=3pt
    },
    tool/.style={
        ellipse, draw=gray!60, fill=gray!8,
        font=\sffamily\tiny, inner sep=2pt, align=center
    },
    chain/.style={->, line width=1.8pt, draw=black!60},
    intflow/.style={->, line width=0.8pt, draw=orange!60, dashed},
    ttp/.style={
        rectangle, rounded corners=2pt, draw=teal!50, fill=teal!5,
        font=\ttfamily\tiny, inner sep=2pt
    },
]

% ── Level nodes (vertical cascade) ───────────────────────────
\node[level=gray]    (L0) at (0, 8)   {L0\\[-2pt]{\footnotesize Script Kiddie}};
\node[level=blue]    (L1) at (0, 6)   {L1\\[-2pt]{\footnotesize Basic}};
\node[level=yellow]  (L2) at (0, 4)   {L2\\[-2pt]{\footnotesize Intermediate}};
\node[level=purple]  (L3) at (0, 2)   {L3\\[-2pt]{\footnotesize Advanced}};
\node[level=red]     (L4) at (0, 0)   {L4\\[-2pt]{\footnotesize APT}};

% ── Cascade arrows ───────────────────────────────────────────
\draw[chain] (L0) -- (L1);
\draw[chain] (L1) -- (L2);
\draw[chain] (L2) -- (L3);
\draw[chain] (L3) -- (L4);

% ── Skill parameter (s) ─────────────────────────────────────
\node[font=\sffamily\tiny\bfseries, text=gray!60] at (-1.9, 8) {$s=0.20$};
\node[font=\sffamily\tiny\bfseries, text=blue!60] at (-1.9, 6) {$s=0.35$};
\node[font=\sffamily\tiny\bfseries, text=yellow!70!black] at (-1.9, 4) {$s=0.55$};
\node[font=\sffamily\tiny\bfseries, text=purple!60] at (-1.9, 2) {$s=0.75$};
\node[font=\sffamily\tiny\bfseries, text=red!60] at (-1.9, 0) {$s=1.00$};

% ── Attack actions (right side) ──────────────────────────────
\node[intel] (A0) at (4.5, 8) {
    \textbullet\ TCP SYN scan (20 ports)\\
    \textbullet\ UDP MAVLink probe\\
    \textbullet\ Banner grabbing
};
\node[intel] (A1) at (4.5, 6) {
    \textbullet\ HEARTBEAT $\to$ sysid 추출\\
    \textbullet\ PARAM\_REQUEST\_LIST\\
    \textbullet\ ARM/TAKEOFF 명령 주입\\
    \textbullet\ GPS spoof 시도
};
\node[intel] (A2) at (4.5, 4) {
    \textbullet\ HTTP API 열거\\
    \textbullet\ 브루트포스 로그인\\
    \textbullet\ 크레덴셜 재사용\\
    \textbullet\ 미션 데이터 추출
};
\node[intel] (A3) at (4.5, 2) {
    \textbullet\ CVE-2026-25253 exploit\\
    \textbullet\ RTSP DoS\\
    \textbullet\ 브레드크럼 체인 추적\\
    \textbullet\ WS skill\_invoke
};
\node[intel] (A4) at (4.5, 0) {
    \textbullet\ SSH lateral movement\\
    \textbullet\ Cross-drone cred reuse\\
    \textbullet\ 펌웨어 업로드\\
    \textbullet\ C2 persistence
};

\draw[->, gray!40] (L0.east) -- (A0.west);
\draw[->, gray!40] (L1.east) -- (A1.west);
\draw[->, gray!40] (L2.east) -- (A2.west);
\draw[->, gray!40] (L3.east) -- (A3.west);
\draw[->, gray!40] (L4.east) -- (A4.west);

% ── Intelligence harvest (far right) ────────────────────────
\node[ttp] (I0) at (9.0, 8) {open\_ports, banners};
\node[ttp] (I1) at (9.0, 6) {sysid, params, arm\_ok};
\node[ttp] (I2) at (9.0, 4) {api\_token, ssh\_pw, signing\_key};
\node[ttp] (I3) at (9.0, 2) {ws\_token, permissions};

\draw[->, gray!30] (A0.east) -- node[above, font=\tiny\sffamily, text=orange!70] {harvest} (I0.west);
\draw[->, gray!30] (A1.east) -- node[above, font=\tiny\sffamily, text=orange!70] {harvest} (I1.west);
\draw[->, gray!30] (A2.east) -- node[above, font=\tiny\sffamily, text=orange!70] {harvest} (I2.west);
\draw[->, gray!30] (A3.east) -- node[above, font=\tiny\sffamily, text=orange!70] {harvest} (I3.west);

% ── Intel chaining arrows (diagonal) ────────────────────────
\draw[intflow] (I0.south) -- (A1.east);
\draw[intflow] (I1.south) -- (A2.east);
\draw[intflow] (I2.south) -- (A3.east);
\draw[intflow] (I3.south) -- (A4.east);

% ── Intel chain label ────────────────────────────────────────
\node[rotate=-90, font=\sffamily\tiny\bfseries, text=orange!60] at (10.3, 4) {Intelligence Chaining};

% ── ATT\&CK TTPs (left side) ────────────────────────────────
\node[ttp] at (-3.2, 8) {T0846 Remote Svc};
\node[ttp] at (-3.2, 6) {T0836 Modify Param};
\node[ttp] at (-3.2, 4) {T0812 Default Cred};
\node[ttp] at (-3.2, 2) {T0821 Modify Task};
\node[ttp] at (-3.2, 0) {T0839 Module FW};

\draw[->, teal!30] (-2.0, 8) -- (L0.west);
\draw[->, teal!30] (-2.0, 6) -- (L1.west);
\draw[->, teal!30] (-2.0, 4) -- (L2.west);
\draw[->, teal!30] (-2.0, 2) -- (L3.west);
\draw[->, teal!30] (-2.0, 0) -- (L4.west);

% ── Protocol labels ──────────────────────────────────────────
\node[font=\sffamily\tiny, fill=white, draw=gray!30, rounded corners=2pt, inner sep=1.5pt] at (2.2, 8) {TCP/UDP};
\node[font=\sffamily\tiny, fill=white, draw=blue!30, rounded corners=2pt, inner sep=1.5pt] at (2.2, 6) {MAVLink};
\node[font=\sffamily\tiny, fill=white, draw=yellow!50, rounded corners=2pt, inner sep=1.5pt] at (2.2, 4) {HTTP};
\node[font=\sffamily\tiny, fill=white, draw=purple!30, rounded corners=2pt, inner sep=1.5pt] at (2.2, 2) {WebSocket};
\node[font=\sffamily\tiny, fill=white, draw=red!30, rounded corners=2pt, inner sep=1.5pt] at (2.2, 0) {SSH/Multi};

\end{tikzpicture}
}
\caption{L0--L4 적응형 공격 모델. 각 레벨은 이전 레벨에서 수확한 인텔리전스를 체이닝하여 다음 레벨의 공격에 활용한다. 오른쪽 점선은 인텔리전스 흐름(Intelligence Chaining)을, 왼쪽은 MITRE ATT\&CK for ICS TTP 매핑을 나타낸다.}
\label{fig:attack}
\end{figure}

% ═══════════════════════════════════════════════════════════════
\begin{figure}[t]
\centering
\resizebox{\columnwidth}{!}{%
\begin{tikzpicture}[
    >=stealth,
    level/.style={
        rectangle, rounded corners=6pt, draw=#1!70, fill=#1!10,
        line width=1pt, minimum width=3.0cm, minimum height=1.2cm,
        align=center, font=\sffamily\bfseries\small
    },
    intel/.style={
        rectangle, rounded corners=3pt, draw=black!40, fill=yellow!8,
        font=\sffamily\scriptsize, align=left, text width=3.5cm,
        inner sep=3pt
    },
    tool/.style={
        ellipse, draw=gray!60, fill=gray!8,
        font=\sffamily\tiny, inner sep=2pt, align=center
    },
    chain/.style={->, line width=1.8pt, draw=black!60},
    intflow/.style={->, line width=0.8pt, draw=orange!60, dashed},
    ttp/.style={
        rectangle, rounded corners=2pt, draw=teal!50, fill=teal!5,
        font=\ttfamily\tiny, inner sep=2pt
    },
]

% ── Level nodes (vertical cascade) ───────────────────────────
\node[level=gray]    (L0) at (0, 8)   {L0\\[-2pt]{\footnotesize Script Kiddie}};
\node[level=blue]    (L1) at (0, 6)   {L1\\[-2pt]{\footnotesize Basic}};
\node[level=yellow]  (L2) at (0, 4)   {L2\\[-2pt]{\footnotesize Intermediate}};
\node[level=purple]  (L3) at (0, 2)   {L3\\[-2pt]{\footnotesize Advanced}};
\node[level=red]     (L4) at (0, 0)   {L4\\[-2pt]{\footnotesize APT}};

% ── Cascade arrows ───────────────────────────────────────────
\draw[chain] (L0) -- (L1);
\draw[chain] (L1) -- (L2);
\draw[chain] (L2) -- (L3);
\draw[chain] (L3) -- (L4);

% ── Skill parameter (s) ─────────────────────────────────────
\node[font=\sffamily\tiny\bfseries, text=gray!60] at (-1.9, 8) {$s=0.20$};
\node[font=\sffamily\tiny\bfseries, text=blue!60] at (-1.9, 6) {$s=0.35$};
\node[font=\sffamily\tiny\bfseries, text=yellow!70!black] at (-1.9, 4) {$s=0.55$};
\node[font=\sffamily\tiny\bfseries, text=purple!60] at (-1.9, 2) {$s=0.75$};
\node[font=\sffamily\tiny\bfseries, text=red!60] at (-1.9, 0) {$s=1.00$};

% ── Attack actions (right side) ──────────────────────────────
\node[intel] (A0) at (4.5, 8) {
    \textbullet\ TCP SYN scan (20 ports)\\
    \textbullet\ UDP MAVLink probe\\
    \textbullet\ Banner grabbing
};
\node[intel] (A1) at (4.5, 6) {
    \textbullet\ HEARTBEAT $\to$ sysid 추출\\
    \textbullet\ PARAM\_REQUEST\_LIST\\
    \textbullet\ ARM/TAKEOFF 명령 주입\\
    \textbullet\ GPS spoof 시도
};
\node[intel] (A2) at (4.5, 4) {
    \textbullet\ HTTP API 열거\\
    \textbullet\ 브루트포스 로그인\\
    \textbullet\ 크레덴셜 재사용\\
    \textbullet\ 미션 데이터 추출
};
\node[intel] (A3) at (4.5, 2) {
    \textbullet\ CVE-2026-25253 exploit\\
    \textbullet\ RTSP DoS\\
    \textbullet\ 브레드크럼 체인 추적\\
    \textbullet\ WS skill\_invoke
};
\node[intel] (A4) at (4.5, 0) {
    \textbullet\ SSH lateral movement\\
    \textbullet\ Cross-drone cred reuse\\
    \textbullet\ 펌웨어 업로드\\
    \textbullet\ C2 persistence
};

\draw[->, gray!40] (L0.east) -- (A0.west);
\draw[->, gray!40] (L1.east) -- (A1.west);
\draw[->, gray!40] (L2.east) -- (A2.west);
\draw[->, gray!40] (L3.east) -- (A3.west);
\draw[->, gray!40] (L4.east) -- (A4.west);

% ── Intelligence harvest (far right) ────────────────────────
\node[ttp] (I0) at (9.0, 8) {open\_ports, banners};
\node[ttp] (I1) at (9.0, 6) {sysid, params, arm\_ok};
\node[ttp] (I2) at (9.0, 4) {api\_token, ssh\_pw, signing\_key};
\node[ttp] (I3) at (9.0, 2) {ws\_token, permissions};

\draw[->, gray!30] (A0.east) -- node[above, font=\tiny\sffamily, text=orange!70] {harvest} (I0.west);
\draw[->, gray!30] (A1.east) -- node[above, font=\tiny\sffamily, text=orange!70] {harvest} (I1.west);
\draw[->, gray!30] (A2.east) -- node[above, font=\tiny\sffamily, text=orange!70] {harvest} (I2.west);
\draw[->, gray!30] (A3.east) -- node[above, font=\tiny\sffamily, text=orange!70] {harvest} (I3.west);

% ── Intel chaining arrows (diagonal) ────────────────────────
\draw[intflow] (I0.south) -- (A1.east);
\draw[intflow] (I1.south) -- (A2.east);
\draw[intflow] (I2.south) -- (A3.east);
\draw[intflow] (I3.south) -- (A4.east);

% ── Intel chain label ────────────────────────────────────────
\node[rotate=-90, font=\sffamily\tiny\bfseries, text=orange!60] at (10.3, 4) {Intelligence Chaining};

% ── ATT\&CK TTPs (left side) ────────────────────────────────
\node[ttp] at (-3.2, 8) {T0846 Remote Svc};
\node[ttp] at (-3.2, 6) {T0836 Modify Param};
\node[ttp] at (-3.2, 4) {T0812 Default Cred};
\node[ttp] at (-3.2, 2) {T0821 Modify Task};
\node[ttp] at (-3.2, 0) {T0839 Module FW};

\draw[->, teal!30] (-2.0, 8) -- (L0.west);
\draw[->, teal!30] (-2.0, 6) -- (L1.west);
\draw[->, teal!30] (-2.0, 4) -- (L2.west);
\draw[->, teal!30] (-2.0, 2) -- (L3.west);
\draw[->, teal!30] (-2.0, 0) -- (L4.west);

% ── Protocol labels ──────────────────────────────────────────
\node[font=\sffamily\tiny, fill=white, draw=gray!30, rounded corners=2pt, inner sep=1.5pt] at (2.2, 8) {TCP/UDP};
\node[font=\sffamily\tiny, fill=white, draw=blue!30, rounded corners=2pt, inner sep=1.5pt] at (2.2, 6) {MAVLink};
\node[font=\sffamily\tiny, fill=white, draw=yellow!50, rounded corners=2pt, inner sep=1.5pt] at (2.2, 4) {HTTP};
\node[font=\sffamily\tiny, fill=white, draw=purple!30, rounded corners=2pt, inner sep=1.5pt] at (2.2, 2) {WebSocket};
\node[font=\sffamily\tiny, fill=white, draw=red!30, rounded corners=2pt, inner sep=1.5pt] at (2.2, 0) {SSH/Multi};

\end{tikzpicture}
}
\caption{L0--L4 적응형 공격 모델. 각 레벨은 이전 레벨에서 수확한 인텔리전스를 체이닝하여 다음 레벨의 공격에 활용한다. 오른쪽 점선은 인텔리전스 흐름(Intelligence Chaining)을, 왼쪽은 MITRE ATT\&CK for ICS TTP 매핑을 나타낸다.}
\label{fig:attack}
\end{figure}

% ═══════════════════════════════════════════════════════════════
\begin{figure}[t]
\centering
\resizebox{\columnwidth}{!}{%
\begin{tikzpicture}[
    >=stealth,
    level/.style={
        rectangle, rounded corners=6pt, draw=#1!70, fill=#1!10,
        line width=1pt, minimum width=3.0cm, minimum height=1.2cm,
        align=center, font=\sffamily\bfseries\small
    },
    intel/.style={
        rectangle, rounded corners=3pt, draw=black!40, fill=yellow!8,
        font=\sffamily\scriptsize, align=left, text width=3.5cm,
        inner sep=3pt
    },
    tool/.style={
        ellipse, draw=gray!60, fill=gray!8,
        font=\sffamily\tiny, inner sep=2pt, align=center
    },
    chain/.style={->, line width=1.8pt, draw=black!60},
    intflow/.style={->, line width=0.8pt, draw=orange!60, dashed},
    ttp/.style={
        rectangle, rounded corners=2pt, draw=teal!50, fill=teal!5,
        font=\ttfamily\tiny, inner sep=2pt
    },
]

% ── Level nodes (vertical cascade) ───────────────────────────
\node[level=gray]    (L0) at (0, 8)   {L0\\[-2pt]{\footnotesize Script Kiddie}};
\node[level=blue]    (L1) at (0, 6)   {L1\\[-2pt]{\footnotesize Basic}};
\node[level=yellow]  (L2) at (0, 4)   {L2\\[-2pt]{\footnotesize Intermediate}};
\node[level=purple]  (L3) at (0, 2)   {L3\\[-2pt]{\footnotesize Advanced}};
\node[level=red]     (L4) at (0, 0)   {L4\\[-2pt]{\footnotesize APT}};

% ── Cascade arrows ───────────────────────────────────────────
\draw[chain] (L0) -- (L1);
\draw[chain] (L1) -- (L2);
\draw[chain] (L2) -- (L3);
\draw[chain] (L3) -- (L4);

% ── Skill parameter (s) ─────────────────────────────────────
\node[font=\sffamily\tiny\bfseries, text=gray!60] at (-1.9, 8) {$s=0.20$};
\node[font=\sffamily\tiny\bfseries, text=blue!60] at (-1.9, 6) {$s=0.35$};
\node[font=\sffamily\tiny\bfseries, text=yellow!70!black] at (-1.9, 4) {$s=0.55$};
\node[font=\sffamily\tiny\bfseries, text=purple!60] at (-1.9, 2) {$s=0.75$};
\node[font=\sffamily\tiny\bfseries, text=red!60] at (-1.9, 0) {$s=1.00$};

% ── Attack actions (right side) ──────────────────────────────
\node[intel] (A0) at (4.5, 8) {
    \textbullet\ TCP SYN scan (20 ports)\\
    \textbullet\ UDP MAVLink probe\\
    \textbullet\ Banner grabbing
};
\node[intel] (A1) at (4.5, 6) {
    \textbullet\ HEARTBEAT $\to$ sysid 추출\\
    \textbullet\ PARAM\_REQUEST\_LIST\\
    \textbullet\ ARM/TAKEOFF 명령 주입\\
    \textbullet\ GPS spoof 시도
};
\node[intel] (A2) at (4.5, 4) {
    \textbullet\ HTTP API 열거\\
    \textbullet\ 브루트포스 로그인\\
    \textbullet\ 크레덴셜 재사용\\
    \textbullet\ 미션 데이터 추출
};
\node[intel] (A3) at (4.5, 2) {
    \textbullet\ CVE-2026-25253 exploit\\
    \textbullet\ RTSP DoS\\
    \textbullet\ 브레드크럼 체인 추적\\
    \textbullet\ WS skill\_invoke
};
\node[intel] (A4) at (4.5, 0) {
    \textbullet\ SSH lateral movement\\
    \textbullet\ Cross-drone cred reuse\\
    \textbullet\ 펌웨어 업로드\\
    \textbullet\ C2 persistence
};

\draw[->, gray!40] (L0.east) -- (A0.west);
\draw[->, gray!40] (L1.east) -- (A1.west);
\draw[->, gray!40] (L2.east) -- (A2.west);
\draw[->, gray!40] (L3.east) -- (A3.west);
\draw[->, gray!40] (L4.east) -- (A4.west);

% ── Intelligence harvest (far right) ────────────────────────
\node[ttp] (I0) at (9.0, 8) {open\_ports, banners};
\node[ttp] (I1) at (9.0, 6) {sysid, params, arm\_ok};
\node[ttp] (I2) at (9.0, 4) {api\_token, ssh\_pw, signing\_key};
\node[ttp] (I3) at (9.0, 2) {ws\_token, permissions};

\draw[->, gray!30] (A0.east) -- node[above, font=\tiny\sffamily, text=orange!70] {harvest} (I0.west);
\draw[->, gray!30] (A1.east) -- node[above, font=\tiny\sffamily, text=orange!70] {harvest} (I1.west);
\draw[->, gray!30] (A2.east) -- node[above, font=\tiny\sffamily, text=orange!70] {harvest} (I2.west);
\draw[->, gray!30] (A3.east) -- node[above, font=\tiny\sffamily, text=orange!70] {harvest} (I3.west);

% ── Intel chaining arrows (diagonal) ────────────────────────
\draw[intflow] (I0.south) -- (A1.east);
\draw[intflow] (I1.south) -- (A2.east);
\draw[intflow] (I2.south) -- (A3.east);
\draw[intflow] (I3.south) -- (A4.east);

% ── Intel chain label ────────────────────────────────────────
\node[rotate=-90, font=\sffamily\tiny\bfseries, text=orange!60] at (10.3, 4) {Intelligence Chaining};

% ── ATT\&CK TTPs (left side) ────────────────────────────────
\node[ttp] at (-3.2, 8) {T0846 Remote Svc};
\node[ttp] at (-3.2, 6) {T0836 Modify Param};
\node[ttp] at (-3.2, 4) {T0812 Default Cred};
\node[ttp] at (-3.2, 2) {T0821 Modify Task};
\node[ttp] at (-3.2, 0) {T0839 Module FW};

\draw[->, teal!30] (-2.0, 8) -- (L0.west);
\draw[->, teal!30] (-2.0, 6) -- (L1.west);
\draw[->, teal!30] (-2.0, 4) -- (L2.west);
\draw[->, teal!30] (-2.0, 2) -- (L3.west);
\draw[->, teal!30] (-2.0, 0) -- (L4.west);

% ── Protocol labels ──────────────────────────────────────────
\node[font=\sffamily\tiny, fill=white, draw=gray!30, rounded corners=2pt, inner sep=1.5pt] at (2.2, 8) {TCP/UDP};
\node[font=\sffamily\tiny, fill=white, draw=blue!30, rounded corners=2pt, inner sep=1.5pt] at (2.2, 6) {MAVLink};
\node[font=\sffamily\tiny, fill=white, draw=yellow!50, rounded corners=2pt, inner sep=1.5pt] at (2.2, 4) {HTTP};
\node[font=\sffamily\tiny, fill=white, draw=purple!30, rounded corners=2pt, inner sep=1.5pt] at (2.2, 2) {WebSocket};
\node[font=\sffamily\tiny, fill=white, draw=red!30, rounded corners=2pt, inner sep=1.5pt] at (2.2, 0) {SSH/Multi};

\end{tikzpicture}
}
\caption{L0--L4 적응형 공격 모델. 각 레벨은 이전 레벨에서 수확한 인텔리전스를 체이닝하여 다음 레벨의 공격에 활용한다. 오른쪽 점선은 인텔리전스 흐름(Intelligence Chaining)을, 왼쪽은 MITRE ATT\&CK for ICS TTP 매핑을 나타낸다.}
\label{fig:attack}
\end{figure}
